\documentclass[../DoAn.tex]{subfiles}
\begin{document}

% Đặt tên caption bảng bằng tiếng Việt khi biên dịch riêng chương này.
\renewcommand{\tablename}{Bảng}

Chương này phân tích các lựa chọn mô hình và thuật toán đã được trình bày ở Chương 2 và Chương 3. Nếu Chương 2 đóng vai trò xây dựng cơ sở lý thuyết, Chương 3 mô tả cách hiện thực hệ thống, thì Chương 4 tập trung trả lời câu hỏi vì sao các lựa chọn đó là hợp lý trong phạm vi đồ án. Nội dung phân tích bao gồm ý nghĩa của chi phí địa hình, không gian trạng thái có hướng, chi phí quay, thuật toán Dijkstra có hướng, chiến lược chọn mục tiêu chưa bao phủ, ràng buộc năng lượng quay về căn cứ, vật cản động, vai trò của mô hình toàn hướng lý tưởng và độ phức tạp tính toán.

\section{Mục tiêu và phạm vi phân tích}

Bài toán bao phủ bản đồ có ràng buộc năng lượng, chi phí quay, vật cản tĩnh và vật cản động là một bài toán có nhiều thành phần phụ thuộc lẫn nhau. Nếu tìm lời giải tối ưu toàn cục cho toàn bộ thứ tự di chuyển của robot, hệ thống phải xét đồng thời vị trí, hướng, năng lượng còn lại, trạng thái bao phủ, trạng thái vật cản động và khả năng quay về căn cứ. Không gian tìm kiếm khi đó tăng rất nhanh theo kích thước bản đồ và số bước thời gian.

Đồ án lựa chọn chiến lược tham lam có ràng buộc vì bài toán được xử lý theo trạng thái hiện tại của robot và môi trường. Ở mỗi lần lập kế hoạch, hệ thống dùng Dijkstra có hướng để tìm đường chi phí nhỏ nhất tới các mục tiêu ứng viên, sau đó chỉ chọn mục tiêu thỏa điều kiện năng lượng và quay về căn cứ. Cách tiếp cận này phù hợp với mô phỏng lặp, cho phép lập kế hoạch lại khi vật cản động thay đổi và giúp log thực nghiệm phản ánh trực tiếp từng quyết định của hệ thống. Mục tiêu phân tích của chương này là làm rõ vì sao lựa chọn đó phù hợp với phạm vi đồ án và với hệ thống đã cài đặt.

Các phân tích trong chương này tập trung vào bốn nhóm câu hỏi chính. Thứ nhất, chi phí của một ô trong bản đồ có ý nghĩa gì và tại sao có thể dùng chi phí chuẩn hóa thay cho mô hình vật lý đầy đủ. Thứ hai, vì sao cần mở rộng trạng thái từ vị trí sang cặp vị trí--hướng. Thứ ba, điều kiện năng lượng quay về căn cứ bảo vệ hệ thống ở mức nào và còn giới hạn gì. Thứ tư, vì sao đồ án chọn Dijkstra có hướng và chiến lược chọn mục tiêu tham lam có ràng buộc thay vì một thuật toán tối ưu toàn cục phức tạp hơn.

\section{Phân tích mô hình chi phí địa hình}

Trong bản đồ lưới có trọng số, mỗi ô hợp lệ $v$ mang một chi phí địa hình $c(v)>0$. Đại lượng này được hiểu là chi phí năng lượng tương đối đã chuẩn hóa của ô đó. Vì vậy, $c(v)$ không nên được hiểu là một nhãn tùy ý, cũng không nhất thiết phải là giá trị năng lượng tuyệt đối theo đơn vị Joule. Ý nghĩa quan trọng nhất của $c(v)$ là quan hệ tương đối giữa các ô: ô khó đi hơn phải có chi phí lớn hơn ô dễ đi hơn trong cùng thang đo mô phỏng.

Cách biểu diễn này phù hợp với mục tiêu lập kế hoạch đường đi. Khi robot xét một bước đi từ ô $u$ sang ô $v$, hệ thống có thể đặt chi phí cạnh bằng chi phí đi vào ô kế tiếp:
\begin{equation}
w(u,v)=c(v).
\end{equation}
Nhờ đó, thuật toán tìm đường không chỉ tối thiểu hóa số bước, mà tối thiểu hóa tổng chi phí địa hình trên đường đi. Nếu một đường có ít bước nhưng đi qua vùng có chi phí lớn, nó có thể bị đánh giá kém hơn một đường dài hơn nhưng đi qua vùng có chi phí thấp. Đây là khác biệt quan trọng giữa tìm đường trên bản đồ không trọng số và tìm đường trên bản đồ có trọng số.

Một điểm cần nhấn mạnh là chi phí trong mô hình có tính chuẩn hóa. Nếu toàn bộ các chi phí địa hình được nhân với cùng một hệ số dương, quan hệ tương đối giữa các ô vẫn được giữ. Khi đó, thứ tự ưu tiên giữa nhiều đường đi thường không thay đổi nếu các thành phần chi phí khác cũng được quy đổi trên cùng thang đo. Điều này cho phép đồ án tập trung vào cấu trúc quyết định của thuật toán thay vì phải mô phỏng đầy đủ mọi đại lượng vật lý của robot ngoài thực địa.

Ví dụ, giả sử có hai đường đi tới cùng một mục tiêu. Đường thứ nhất chỉ gồm 5 bước nhưng đi qua các ô có chi phí cao, tổng chi phí là 20. Đường thứ hai gồm 8 bước nhưng đi qua nền dễ di chuyển, tổng chi phí là 8. Nếu chỉ xét số bước, đường thứ nhất có vẻ tốt hơn. Nếu xét năng lượng chuẩn hóa, đường thứ hai hợp lý hơn. Đây là lý do việc đưa $c(v)$ vào mô hình không chỉ làm bản đồ nhiều thông tin hơn, mà trực tiếp thay đổi quyết định lập kế hoạch của robot.

Với vật cản tĩnh hoặc ô không an toàn, hệ thống có thể gán chi phí hiệu dụng bằng $\infty$:
\begin{equation}
c_{\text{eff}}(v)=\infty.
\end{equation}
Về mặt thuật toán, thao tác này tương đương với việc loại ô đó khỏi đồ thị tìm đường. Cách biểu diễn bằng $\infty$ thuận tiện cho cài đặt vì hệ thống vẫn giữ cấu trúc ma trận bản đồ, nhưng có thể thay đổi nhanh trạng thái đi được hoặc không đi được của một ô theo thời gian.

\section{Phân tích không gian trạng thái có hướng}

Nếu robot chỉ được biểu diễn bằng vị trí $v$, thuật toán không phân biệt được hai trường hợp robot đứng cùng một ô nhưng quay về hai hướng khác nhau. Điều này không gây vấn đề lớn khi chi phí quay bằng 0. Tuy nhiên, khi thao tác đổi hướng tiêu tốn thời gian hoặc năng lượng, hai trạng thái cùng vị trí nhưng khác hướng không còn tương đương.

Vì vậy, đồ án mở rộng trạng thái tìm đường thành:
\begin{equation}
q=(v,h),
\end{equation}
trong đó $v$ là ô hiện tại và $h$ là hướng hiện tại của robot. Với mô hình lưới bốn hướng:
\begin{equation}
H=\{N,E,S,W\}.
\end{equation}
Không gian trạng thái có hướng là:
\begin{equation}
Q=V_{\text{free}}\times H.
\end{equation}
Gọi $N=|V_{\text{free}}|$, ta có:
\begin{equation}
|Q|=4N.
\end{equation}
Như vậy, so với mô hình chỉ xét vị trí, số trạng thái tăng bốn lần. Đây là chi phí tính toán phải trả để thuật toán có khả năng xét chi phí quay. Đổi lại, hệ thống có thể phân biệt được tình huống robot đang hướng Đông và đi tiếp sang ô phía Đông với tình huống robot đang hướng Bắc nhưng muốn đi sang ô phía Đông. Hai thao tác này có cùng đích đến, nhưng khác chi phí do số lần quay khác nhau.

Mô hình bốn hướng là lựa chọn trực tiếp cho bản đồ lưới vuông. Nếu xét chuyển động chéo, chuyển động tám hướng hoặc góc quay mịn hơn, có thể mở rộng tập $H$ và định nghĩa thêm các cạnh chuyển trạng thái tương ứng. Bản chất bài toán không thay đổi: robot vẫn tìm đường trên không gian trạng thái có hướng. Tuy nhiên, số trạng thái, số cạnh và số trường hợp chuyển hướng đều tăng lên. Vì vậy, đồ án giữ bốn hướng để giới hạn độ phức tạp và bám sát cách cài đặt của chương trình mô phỏng.

\section{Phân tích chi phí quay và lựa chọn chuẩn hóa}

Chi phí chuyển từ trạng thái $(v,h)$ sang trạng thái $(v',h')$ trong hệ thống có dạng:
\begin{equation}
w_q\bigl((v,h),(v',h')\bigr)=c_{\text{eff}}(v')+\Delta(h,h')\tau(v),
\end{equation}
trong đó $c_{\text{eff}}(v')$ là chi phí đi vào ô kế tiếp, $\Delta(h,h')$ là số lần quay 90 độ nhỏ nhất giữa hai hướng, còn $\tau(v)$ là chi phí cho một lần quay 90 độ tại ô hiện tại.

Trong mô hình mặc định của đồ án:
\begin{equation}
\tau(v)=\frac{1}{2}c(v).
\end{equation}
Do đó, nếu robot đi thẳng thì $\Delta(h,h')=0$ và chi phí chỉ là chi phí đi vào ô kế tiếp. Nếu robot quay 90 độ trước khi đi, chi phí tăng thêm $\frac{1}{2}c(v)$. Nếu robot quay 180 độ, chi phí tăng thêm $c(v)$.

Cách chọn này xuất phát từ chuẩn hóa đã nêu ở Chương 2. Với robot vi sai hai bánh, nếu gọi $L$ là khoảng cách giữa hai bánh và $s$ là độ dài cạnh của một ô lưới, chi phí quay tại chỗ có thể được mô hình hóa tỷ lệ với:
\begin{equation}
\frac{L|\theta|}{2s}.
\end{equation}
Khi chọn:
\begin{equation}
L=\frac{2s}{\pi},
\end{equation}
ta có:
\begin{equation}
E_{\text{turn}}(\theta,v)=\kappa_{\text{turn}}c(v)\frac{|\theta|}{\pi}.
\end{equation}
Với $\theta=\frac{\pi}{2}$ và $\kappa_{\text{turn}}=1$, chi phí quay 90 độ bằng:
\begin{equation}
E_{\text{turn}}\left(\frac{\pi}{2},v\right)=\frac{1}{2}c(v).
\end{equation}
Với $\theta=\pi$, chi phí quay 180 độ bằng:
\begin{equation}
E_{\text{turn}}(\pi,v)=c(v).
\end{equation}
Bước lượng tử năng lượng $0.5$ trong triển khai cũng xuất phát từ lựa chọn này. Vì chi phí quay 90 độ trong mô hình mặc định bằng một nửa chi phí địa hình tại ô hiện tại, hệ thống cần một đơn vị ghi nhận nhỏ hơn 1 để không làm mất phần chi phí quay này. Lượng tử hóa theo bước $0.5$ giúp log và thống kê năng lượng dễ đọc, đồng thời vẫn giữ được khác biệt giữa đi thẳng, quay 90 độ và quay 180 độ.

Điểm quan trọng là lựa chọn $L=2s/\pi$ không phải là khẳng định rằng mọi robot vật lý đều có kích thước như vậy. Đây là một chuẩn hóa để đưa chi phí quay và chi phí di chuyển về cùng một thang đo dễ diễn giải. Nếu robot có tỷ lệ $L/s$ khác, chi phí quay vẫn thay đổi tỉ lệ thuận theo $L/s$. Sự khác biệt đó có thể được biểu diễn bằng cách thay đổi hệ số $\kappa_{\text{turn}}$ hoặc hệ số quay tương ứng. Nói cách khác, đồ án chọn một trường hợp tiêu biểu để cài đặt và đánh giá, còn các tỷ lệ hình học khác có thể quy đổi trong cùng mô hình.

Cách phân tách này cũng giúp mô hình bao quát nhiều cấu trúc robot khác nhau. Với robot vi sai hai bánh hoặc robot có thân và hướng di chuyển rõ ràng, chi phí quay có thể đáng kể và cần được xét trong lập kế hoạch. Với robot toàn hướng dùng bánh omni, hoặc các hệ có khả năng đổi hướng linh hoạt, chi phí quay có thể nhỏ không đáng kể trong phạm vi mô hình lập kế hoạch. Khi đó, mô hình toàn hướng lý tưởng đóng vai trò đối chứng bằng cách đặt:
\begin{equation}
\tau(v)=0.
\end{equation}
Như vậy, nhiều cấu trúc truyền động khác nhau có thể được nhìn như các trường hợp khác nhau của cùng một mô hình chi phí quay.

\section{Phân tích Dijkstra có hướng}

Dijkstra có hướng trong đồ án thực chất là Dijkstra trên đồ thị trạng thái $Q=V_{\text{free}}\times H$. Mỗi trạng thái là một cặp vị trí--hướng. Từ một trạng thái, thuật toán xét tối đa bốn hướng di chuyển kế tiếp. Nếu ô kế tiếp hợp lệ và an toàn, thuật toán tính chi phí chuyển trạng thái rồi thực hiện bước cập nhật:
\begin{equation}
d(q') \leftarrow \min\bigl(d(q'), d(q)+w_q(q,q')\bigr).
\end{equation}

Vì các chi phí địa hình dương và chi phí quay không âm, trọng số cạnh trong đồ thị trạng thái là không âm. Điều này phù hợp với điều kiện sử dụng Dijkstra. Do đó, tại một thời điểm lập kế hoạch cố định, Dijkstra có hướng tìm được đường có chi phí nhỏ nhất từ trạng thái hiện tại tới các trạng thái đích theo mô hình chi phí đang xét.

Tuy nhiên, cần phân biệt rõ hai mức tối ưu. Dijkstra có hướng tối ưu cho bài toán tìm đường giữa một trạng thái bắt đầu và một mục tiêu cụ thể tại thời điểm lập kế hoạch. Nó không tự động tạo ra thứ tự bao phủ tối ưu cho toàn bộ nhiệm vụ. Bài toán bao phủ toàn cục còn liên quan tới việc chọn mục tiêu nào trước, quay về căn cứ khi nào, xử lý vật cản động ra sao và sạc lại bao nhiêu lần. Vì vậy, Dijkstra là thành phần tìm đường cục bộ tối ưu, còn toàn bộ hệ thống vẫn là một chính sách bao phủ có ràng buộc.

Ưu điểm chính của Dijkstra có hướng trong đồ án là mọi thành phần chi phí đều đi trực tiếp vào trọng số cạnh: chi phí đi vào ô kế tiếp, chi phí quay tại ô hiện tại và trạng thái an toàn của ô. Vì vậy, khi thay đổi mô hình chi phí, đường đi được chọn cũng có thể thay đổi, chứ không chỉ tổng năng lượng sau cùng thay đổi.

Một lựa chọn triển khai liên quan là hướng ban đầu của robot. Nếu hướng ban đầu được đặt tùy ý, mô hình có thể tạo ra chi phí quay phụ thuộc vào lựa chọn khởi tạo hơn là phụ thuộc vào cấu trúc bản đồ. Vì vậy, hệ thống đánh giá cả bốn hướng ban đầu bằng cùng mô hình Dijkstra có hướng. Tiêu chí chọn hướng gồm số mục tiêu có thể tiếp cận, chi phí trung bình tới các mục tiêu đó và độ dài vùng trống phía trước. Cách làm này không nhằm tối ưu toàn bộ nhiệm vụ từ thời điểm ban đầu, mà nhằm chọn một hướng xuất phát hợp lý, không dựa trên một giới hạn mẫu tùy ý và không cần thêm tham số khó giải thích.

\section{Phân tích chiến lược chọn mục tiêu bao phủ}

Tại mỗi thời điểm cần lập kế hoạch, hệ thống xét tập các ô chưa bao phủ:
\begin{equation}
U_t=\{v\in V_{\text{free}}\mid \text{covered}_t(v)=false\}.
\end{equation}
Với mỗi ứng viên $g\in U_t$, hệ thống ước lượng chi phí đi từ vị trí hiện tại tới mục tiêu:
\begin{equation}
d_{\text{to}}(g)=E(P_{v_t\rightarrow g}),
\end{equation}
và chi phí từ mục tiêu quay về căn cứ:
\begin{equation}
d_{\text{return}}(g)=E(P_{g\rightarrow b}).
\end{equation}
Một mục tiêu chỉ được xem là khả thi nếu:
\begin{equation}
E_t \ge d_{\text{to}}(g)+d_{\text{return}}(g)+M(d_{\text{return}}(g)).
\end{equation}
Trong đó $M(d_{\text{return}}(g))$ là phần năng lượng dự phòng cho đường quay về.

Chiến lược chọn mục tiêu của đồ án có thể xem là chiến lược tham lam có ràng buộc năng lượng. Robot ưu tiên các mục tiêu có chi phí đi tới thấp, nhưng chỉ chọn mục tiêu nếu mục tiêu đó thỏa mãn điều kiện quay về căn cứ. Như vậy, hệ thống không tham lam tuyệt đối theo mục tiêu gần nhất, mà là tham lam trong tập các mục tiêu khả thi.

Cách chọn này có ba ưu điểm thực tế. Thứ nhất, nó đơn giản và chạy được trong vòng lặp mô phỏng. Thứ hai, nó kết hợp trực tiếp với Dijkstra có hướng và mô hình chi phí địa hình. Thứ ba, khi vật cản động làm ảnh hưởng đường hiện tại, hệ thống có thể lập kế hoạch lại từ đúng trạng thái hiện tại của robot.

Đổi lại, chiến lược này không bảo đảm thứ tự bao phủ tối ưu toàn cục. Có những trường hợp một mục tiêu gần tại thời điểm hiện tại không phải là lựa chọn tốt nhất nếu xét toàn bộ chuỗi mục tiêu tương lai. Các hướng mở rộng như tối ưu thứ tự thăm, biến thể TSP/VRP, quy hoạch nguyên hoặc học tăng cường có thể cho lời giải mạnh hơn, nhưng cũng làm tăng mạnh độ phức tạp tính toán và độ khó kiểm chứng. Trong phạm vi đồ án, chiến lược tham lam có ràng buộc là lựa chọn phù hợp vì chạy được trong mô phỏng, giải thích được bằng log và bám sát các ràng buộc năng lượng.

\section{Phân tích ràng buộc năng lượng quay về căn cứ}

Ràng buộc quay về căn cứ là một điểm quan trọng để tránh việc robot đi quá xa rồi không còn đủ năng lượng để quay lại. Điều kiện:
\begin{equation}
E_t \ge d_{\text{to}}(g)+d_{\text{return}}(g)+M(d_{\text{return}}(g))
\end{equation}
có nghĩa là robot chỉ chọn mục tiêu $g$ nếu năng lượng hiện tại đủ cho ba phần: đi tới mục tiêu, từ mục tiêu quay về căn cứ và giữ lại một phần dự phòng.

Điều kiện này là một bảo đảm có điều kiện trong mô hình tại thời điểm lập kế hoạch. Nó không phải là bảo đảm tuyệt đối cho mọi diễn biến tương lai, vì vật cản động có thể xuất hiện, đường về có thể bị chặn tạm thời hoặc robot phải chờ và lập kế hoạch lại. Vì vậy, hệ thống không chỉ kiểm tra năng lượng một lần, mà còn có trạng thái quay về căn cứ, trạng thái cảnh báo, trạng thái chờ an toàn và cơ chế tìm đường lại khi cần.

Trong triển khai hiện tại, phần dự phòng được áp dụng theo loại kịch bản:
\begin{equation}
M(d)=
\begin{cases}
0, & \text{khi kịch bản không có vật cản động},\\
\max(M_{\min}, M_{\text{prop}}(d)), & \text{khi kịch bản có vật cản động}.
\end{cases}
\end{equation}
Cách thiết kế này tách rõ hai tình huống. Với bản đồ tĩnh, đường quay về được đánh giá trên bản đồ đã biết; vì vậy robot chỉ cần đủ năng lượng để đi tới mục tiêu và quay về căn cứ. Với kịch bản có vật cản động, đường về ban đầu có thể không còn khả dụng tại thời điểm robot quay về, nên hệ thống giữ thêm quỹ dự phòng.

Trong nhánh có vật cản động, thành phần tỷ lệ được chọn xấp xỉ bằng $25\%$ chi phí đường quay về, tương ứng với hệ số chia 4:
\begin{equation}
M_{\text{prop}}(d)\approx \frac{d}{4}.
\end{equation}
Mức $25\%$ được hiểu là quỹ dự phòng cho chi phí phát sinh khi quay về căn cứ. Nếu vật cản động chặn đường về dự kiến, robot có thể phải chờ, đổi hướng, tìm đường vòng hoặc lập kế hoạch lại; các hành vi này đều có thể làm chi phí thực tế lớn hơn chi phí ước lượng ban đầu. Ngược lại, trong kịch bản tĩnh, giữ cùng một quỹ dự phòng sẽ làm robot quay về sớm hơn dù không có rủi ro động tương ứng.

Bên cạnh thành phần tỷ lệ, hệ thống dùng mức dự phòng tối thiểu $M_{\min}=15$ trong các kịch bản có vật cản động. Giá trị này bảo đảm rằng ngay cả khi đường quay về ngắn, robot vẫn giữ lại một khoảng đệm năng lượng cho các thao tác phát sinh như quay, chờ, lập kế hoạch lại hoặc đi vòng ngắn. Nếu chỉ dùng tỷ lệ theo $d$, các đường về rất ngắn sẽ tạo ra mức dự phòng quá nhỏ.

Dạng này không phải công thức dự phòng duy nhất. Có thể thiết kế dự phòng cố định, dự phòng theo phần trăm dung lượng pin, dự phòng theo xác suất gặp vật cản hoặc dự phòng theo mức rủi ro địa hình. Tuy nhiên, điểm quan trọng của phiên bản hiện tại là dự phòng được gắn với rủi ro môi trường, thay vì áp dụng máy móc cho mọi kịch bản. Điều này giúp chính sách năng lượng phản ánh đúng mục tiêu của quỹ dự phòng: bảo vệ khả năng quay về trước rủi ro thay đổi môi trường.

Ngưỡng năng lượng tới hạn $E_{\min}=3$ được dùng như tín hiệu cho thấy robot chỉ còn rất ít năng lượng để thực hiện thêm thao tác trong mô hình. Đây không phải là ngưỡng thất bại tuyệt đối, mà là điều kiện để hệ thống chuyển sang hành vi thận trọng hơn như ưu tiên quay về, chờ chỉ thị hoặc bảo toàn trạng thái nhiệm vụ.

Một chi tiết thực tế khác là robot về tới căn cứ với đúng 0 năng lượng vẫn được xem là quay về thành công, vì robot đã tới điểm sạc hoặc điểm kết thúc nhiệm vụ. Do đó, khi kiểm tra đường về, hệ thống chấp nhận trường hợp chi phí đường về bằng đúng năng lượng còn lại nếu ô đích cuối cùng là căn cứ; chỉ khi chi phí vượt quá năng lượng còn lại thì đường về mới bị loại.

Trong phiên bản hiện tại, đồ án sử dụng một căn cứ duy nhất trùng với vị trí xuất phát. Đây là lựa chọn đơn giản để tập trung vào cơ chế quay về và sạc. Nếu mở rộng sang nhiều căn cứ hoặc nhiều trạm sạc, có thể thay căn cứ đơn $b$ bằng tập căn cứ $B$ và định nghĩa:
\begin{equation}
d_{\text{return}}(g)=\min_{b\in B}E(P_{g\rightarrow b}).
\end{equation}
Khi đó bản chất của điều kiện năng lượng không thay đổi, chỉ thay đổi cách tính đường quay về khả thi nhất.

\section{Phân tích vật cản động và cơ chế lập kế hoạch lại}

Trong môi trường tĩnh, một đường đi đã được tính toán có thể được xem là hợp lệ cho tới khi robot đi hết đường đó. Trong môi trường có vật cản động, giả định này không còn đúng. Một ô an toàn tại thời điểm $t$ có thể trở thành không an toàn tại thời điểm $t+1$. Vì vậy, hệ thống phải xem đường đi hiện tại là một kế hoạch tạm thời, không phải cam kết bất biến.

Gọi $O_t$ là tập ô không an toàn tại thời điểm $t$, tập ô an toàn được định nghĩa là:
\begin{equation}
S_t=V_{\text{free}}\setminus O_t.
\end{equation}
Một bước đi tới ô $v'$ chỉ hợp lệ nếu:
\begin{equation}
v'\in S_t.
\end{equation}
Khi phát hiện ô kế tiếp không an toàn, robot không tiếp tục thực hiện đường đi cũ. Thay vào đó, hệ thống có thể chờ vật cản rời đi, chuyển sang trạng thái cảnh báo, lập kế hoạch lại hoặc quay về căn cứ tùy theo trạng thái năng lượng và nhiệm vụ.

Điểm đáng chú ý là vật cản động không làm Dijkstra có hướng sai về mặt thuật toán. Dijkstra vẫn tìm đường tốt nhất trên trạng thái bản đồ mà nó nhận tại thời điểm lập kế hoạch. Vấn đề là trạng thái bản đồ thay đổi theo thời gian. Vì vậy, giải pháp phù hợp là kết hợp Dijkstra với cơ chế kiểm tra an toàn liên tục và lập kế hoạch lại khi đường hiện tại không còn hợp lệ.

Tập $O_t$ cũng có thể được hiểu rộng hơn là vùng không an toàn, không nhất thiết chỉ gồm đúng ô mà vật cản đang chiếm. Nếu muốn xét kích thước vật cản lớn hơn một ô, khoảng cách an toàn, vùng đệm quanh người hoặc sai số định vị, có thể giãn nở $O_t$. Thuật toán không thay đổi về bản chất, vì điều kiện cốt lõi vẫn là kiểm tra ô kế tiếp có thuộc $S_t$ hay không.

\section{Phân tích mô hình toàn hướng lý tưởng như mô hình đối chứng}

Trong các file log và nhóm kịch bản thực nghiệm, cấu hình này thường được nhận diện bằng tên Metric H hoặc nhóm H. Trong phần phân tích lý thuyết, báo cáo dùng tên chính là mô hình toàn hướng lý tưởng để nhấn mạnh giả định cơ học và tránh biến tên cấu hình thành thuật ngữ học thuật chính.

Mô hình toàn hướng lý tưởng là mô hình đối chứng trong đó chi phí quay được bỏ qua ở mức lập kế hoạch:
\begin{equation}
\tau(v)=0.
\end{equation}
Khi đó, chi phí chuyển trạng thái chỉ còn:
\begin{equation}
w_q\bigl((v,h),(v',h')\bigr)=c_{\text{eff}}(v').
\end{equation}
Mô hình toàn hướng lý tưởng là mô hình đối chứng, không phải mô hình chính của đồ án. Nó biểu diễn giả định rằng thao tác đổi hướng có chi phí nhỏ không đáng kể trong phạm vi lập kế hoạch. Giả định này phù hợp với các hệ có khả năng đổi hướng linh hoạt hoặc khi cần tách riêng ảnh hưởng của chi phí quay.

\begin{table}[!htbp]
\centering
\caption{So sánh mô hình có chi phí quay và mô hình toàn hướng}
\small
\setlength{\tabcolsep}{4pt}
\renewcommand{\arraystretch}{1.05}
\begin{tabular}{|L{0.25\textwidth}|L{0.32\textwidth}|L{0.33\textwidth}|}
\hline
\textbf{Tiêu chí} & \textbf{Mô hình có chi phí quay} & \textbf{Mô hình toàn hướng lý tưởng} \\
\hline
Chi phí quay & Được tính vào trọng số cạnh thông qua $\tau(v)$. & Được bỏ qua ở mức lập kế hoạch, $\tau(v)=0$. \\
\hline
Ảnh hưởng của hướng & Hướng hiện tại ảnh hưởng trực tiếp tới chi phí bước kế tiếp. & Hướng còn được lưu trong trạng thái nhưng không làm tăng chi phí quay. \\
\hline
Xu hướng đường đi & Có thể ưu tiên đường ít đổi hướng hơn nếu tổng chi phí thấp hơn. & Có thể đổi hướng nhiều hơn vì thao tác quay không bị phạt năng lượng. \\
\hline
Vai trò trong đồ án & Mô hình chính để khảo sát robot có thân và hướng chuyển động rõ ràng. & Mô hình đối chứng để đánh giá ảnh hưởng của việc bỏ qua chi phí quay. \\
\hline
\end{tabular}\end{table}

Việc so sánh mô hình toàn hướng lý tưởng, tương ứng với cấu hình Metric H trong thực nghiệm, với mô hình có chi phí quay giúp trả lời một câu hỏi quan trọng: chi phí quay chỉ làm tăng tổng năng lượng sau khi robot đã chọn đường, hay nó làm thay đổi cả đường đi được chọn. Do chi phí quay nằm trực tiếp trong trọng số cạnh của Dijkstra có hướng, nó có thể ảnh hưởng ngay trong quá trình lập kế hoạch. Vì vậy, nếu hai mô hình cho hành vi khác nhau về số bước, số lần quay về, số ô đi lại hoặc năng lượng tiêu thụ, sự khác biệt đó phản ánh tác động thực sự của giả định chi phí quay.

\section{Phân tích độ phức tạp tính toán}

Gọi $N=|V_{\text{free}}|$ là số ô hợp lệ trong bản đồ. Với bốn hướng, số trạng thái có hướng là:
\begin{equation}
|Q|=4N.
\end{equation}
Từ mỗi trạng thái, robot xét tối đa bốn hướng di chuyển kế tiếp. Do đó số cạnh trạng thái có thể được chặn trên bởi:
\begin{equation}
|E_Q|\le 16N.
\end{equation}
Dijkstra dùng hàng đợi ưu tiên có độ phức tạp:
\begin{equation}
O\bigl((|Q|+|E_Q|)\log |Q|\bigr).
\end{equation}
Thay các chặn trên vào, ta có:
\begin{equation}
O(N\log N),
\end{equation}
vì các hệ số 4 và 16 là hằng số đối với mô hình bốn hướng.

Khi chọn mục tiêu bao phủ, hệ thống không chỉ chạy một lần tìm đường. Trước hết, Dijkstra được dùng để tính chi phí từ robot tới các ô chưa bao phủ. Sau đó, với một số ứng viên, hệ thống cần ước lượng chi phí từ mục tiêu đó quay về căn cứ. Nếu trong một lần chọn mục tiêu có $k$ ứng viên được đánh giá sâu bằng Dijkstra có hướng, chi phí có thể mô tả ở mức khái quát là:
\begin{equation}
O\bigl((1+k)N\log N\bigr),
\end{equation}
chưa tính chi phí sắp xếp danh sách ứng viên. Nếu số ứng viên là $r$, thao tác sắp xếp có thể tốn:
\begin{equation}
O(r\log r).
\end{equation}
Trong trường hợp xấu, $r$ có thể cùng bậc với $N$.

\begin{table}[!htbp]
\centering
\caption{Độ phức tạp khái quát của các thành phần lập kế hoạch}
\small
\setlength{\tabcolsep}{4pt}
\renewcommand{\arraystretch}{1.05}
\begin{tabular}{|L{0.34\textwidth}|L{0.24\textwidth}|L{0.32\textwidth}|}
\hline
\textbf{Thành phần xử lý} & \textbf{Độ phức tạp khái quát} & \textbf{Ghi chú} \\
\hline
Một lần Dijkstra có hướng & $O(N\log N)$ & Với $N$ là số ô hợp lệ và bốn hướng cố định. \\
\hline
Thu thập ứng viên chưa bao phủ & $O(N)$ & Duyệt các ô hợp lệ và lấy các ô chưa bao phủ còn tới được. \\
\hline
Sắp xếp ứng viên & $O(r\log r)$ & $r$ là số ứng viên được tìm thấy. \\
\hline
Kiểm tra $k$ ứng viên sâu & $O(kN\log N)$ & Mỗi ứng viên có thể cần ước lượng đường quay về căn cứ. \\
\hline
Lập kế hoạch lại nhiều lần & Phụ thuộc số lần replan & Môi trường động có thể làm tăng số lần gọi Dijkstra. \\
\hline
\end{tabular}\end{table}

Phân tích trên cho thấy việc chọn Dijkstra có hướng và chiến lược tham lam có ràng buộc là lựa chọn thực dụng. Hệ thống có thể chạy lặp lại nhiều lần trên bản đồ mô phỏng, trong khi vẫn giữ được cơ sở tối ưu cục bộ cho từng lần tìm đường. Nếu thay bằng tối ưu toàn cục trên toàn bộ thứ tự bao phủ, không gian tìm kiếm sẽ lớn hơn nhiều và khó kiểm soát hơn trong phạm vi đồ án.

\section{Giới hạn của phương pháp và khả năng mở rộng}

Các giới hạn của phương pháp xác định ranh giới của phiên bản mô phỏng hiện tại. Đồ án tập trung vào lập kế hoạch bao phủ trên bản đồ lưới đã biết, vì đây là phần trực tiếp gắn với thuật toán, trạng thái có hướng, năng lượng và cơ chế quay về căn cứ. Các bài toán như xây dựng bản đồ, định vị, điều khiển liên tục, mô hình cơ học đất hoặc cảm biến dò mìn thuộc các lớp bài toán khác và có thể được tích hợp ở các hướng phát triển sau.

Trong phạm vi đó, các lựa chọn của đồ án tạo thành một khung mô phỏng cụ thể và kiểm chứng được. Một căn cứ có thể mở rộng thành nhiều căn cứ; vật cản động một ô có thể mở rộng thành vùng không an toàn; bốn hướng có thể mở rộng thành nhiều hướng; hệ số chi phí quay có thể hiệu chỉnh theo loại robot; và chính sách dự phòng $M(d)$ có thể thay bằng hàm phụ thuộc mức rủi ro môi trường.

\begin{table}[!htbp]
\centering
\caption{Các lựa chọn đại diện và khả năng mở rộng của mô hình}
\small
\setlength{\tabcolsep}{4pt}
\renewcommand{\arraystretch}{1.05}
\begin{tabular}{|L{0.30\textwidth}|L{0.30\textwidth}|L{0.30\textwidth}|}
\hline
\textbf{Lựa chọn trong đồ án} & \textbf{Ý nghĩa} & \textbf{Khả năng mở rộng} \\
\hline
Bản đồ lưới bốn hướng & Giữ mô hình đơn giản và rõ ràng. & Mở rộng tập hướng $H$ hoặc thêm cạnh chéo. \\
\hline
Một căn cứ $b$ & Tập trung vào quay về và sạc. & Mở rộng thành tập căn cứ $B$. \\
\hline
Chi phí quay $\tau(v)=\frac{1}{2}c(v)$ & Chuẩn hóa chi phí quay 90 độ. & Thay đổi $\kappa_{\text{turn}}$ hoặc hệ số theo loại robot. \\
\hline
Vật cản động dạng ô & Dễ cập nhật trạng thái an toàn. & Giãn nở thành vùng không an toàn $O_t$. \\
\hline
Tham lam có ràng buộc năng lượng & Chạy được trong mô phỏng lặp và bám theo trạng thái hiện tại. & Thay bằng tối ưu toàn cục hoặc chiến lược học nếu mở rộng nghiên cứu. \\
\hline
\end{tabular}\end{table}

Các lựa chọn trên xác định phiên bản hiện tại của hệ thống mô phỏng. Chúng đủ cụ thể để cài đặt và đánh giá, đồng thời vẫn mở đường cho các biến thể sau này như nhiều căn cứ, nhiều hướng di chuyển hơn hoặc hàm dự phòng phụ thuộc rủi ro.

\section{Kết chương}

Chương này đã phân tích các lựa chọn lý thuyết và thiết kế chính của đồ án. Chi phí địa hình $c(v)$ được đưa vào đồ thị có trọng số để robot chọn đường theo năng lượng thay vì chỉ theo số bước. Việc mở rộng trạng thái thành $(v,h)$ đưa chi phí quay trực tiếp vào quá trình lập kế hoạch. Dijkstra có hướng tìm đường chi phí nhỏ nhất giữa các trạng thái trong mô hình tại thời điểm lập kế hoạch; phần chọn mục tiêu và quay về căn cứ được xử lý bằng chính sách tham lam có ràng buộc.

Chương này cũng đã phân tích chiến lược chọn mục tiêu tham lam có ràng buộc năng lượng, điều kiện quay về căn cứ, vai trò của hàm dự phòng có điều kiện $M(d)$, cơ chế xử lý vật cản động thông qua tập ô an toàn, vai trò đối chứng của mô hình toàn hướng lý tưởng và độ phức tạp tính toán của các thành phần chính. Các phân tích này tạo cơ sở để Chương 5 đánh giá thực nghiệm hành vi của hệ thống trên các bản đồ và cấu hình mô phỏng cụ thể.

\end{document}
